% Source: sources/returns/S6_04_MAPS_MANIFOLDS/sections/maps_weighted_engel_invariants.tex:3–301
\paragraph{Theorem D.2.48 (Proved: complete smooth classification of the simply connected weighted Engel family).}
Take the weight-$d$ Engel construction of D.2.34, for any $d>0$,
with both additive sites replaced by good logarithmic caps.  Retain
all the actual marked data
\[
 m_0=3k,\qquad m_1=4\ell,\qquad
 v_i=r_i\delta+n_i u_i,\qquad
 P=m_1n_0+m_0n_1=\pm1 .
 \tag{D.WI1}
\]
The manifold is simply connected.  Its complex first and second
Chern classes, spin characteristic class, and refined class satisfy
\[
 c_1(X)=c_2(X)=\lambda(X)=0,\qquad P_{\mathrm{ref}}(X)=0.
 \tag{D.WI2}
\]
Here $P_{\mathrm{ref}}$ denotes Zhubr's class, to distinguish it from
the Seifert determinant $P$ in (D.WI1).

The classes $x_i=\operatorname{PD}(D_i)$ of the $d$ cusp components,
indexed cyclically, generate $H^2(X;\mathbb Z)$ with sole relation
$\sum_i x_i=0$.  The complex ruling $f=H-E_1$ has image
$s=[a\wedge b]$, a generator of $\operatorname{Tor}H_2(X)=\mathbb Z/d$.
Put $\kappa=\operatorname{PD}(f)$.  With repetitions retained when
$d=2$, the complete degree-two product is
\[
 x_i x_j=
 \bigl(\mathbf1_{j=i+1}+\mathbf1_{j=i-1}-2\mathbf1_{j=i}\bigr)\kappa.
 \tag{D.WI3}
\]
For $d=1$ the sole component class is zero.  For $d=2$ the two
off-diagonal incidences both occur, and every product in (D.WI3)
is zero since $2\kappa=0$.  The integral cubic is zero for every $d$.
Here is the complete fundamental-class invariant, including its
two- and three-primary coordinates. Put
$r=d-1$, $D_2=d\gcd(d,2)$, and $D_3=d\gcd(d,3)$.
For $d>1$ there is a class $u\in H^2(X;\mathbb Z/d)$
which, together with $x_0,\ldots,x_{d-2}$, fixes the splitting
$H_2(X)=\mathbb Z^r\oplus\mathbb Z/d$ described below. In
\[
 H_6(K(\mathbb Z^r\oplus\mathbb Z/d,2);\mathbb Z)
 \simeq
 \mathbb Z^{\binom{r+2}{3}}\oplus
 (\mathbb Z/d)^{\binom{r+1}{2}}\oplus
 (\mathbb Z/D_2)^r\oplus\mathbb Z/D_3
 \tag{D.WI14}
\]
the complete class $\mu$ is determined by
\[
 \begin{aligned}
 \langle x_i x_j x_k,[X]\rangle&=0,\\
 \langle u x_i x_j,[X]\rangle
    &=\mathbf1_{j=i+1}+\mathbf1_{j=i-1}-2\mathbf1_{j=i}
                                                     &&\pmod d,\\
 \langle \mathcal Q_2(u)x_i,[X]\rangle&=-12             &&\pmod {D_2},\\
 \langle \mathcal Q_3(u),[X]\rangle&=0                 &&\pmod {D_3}.
 \end{aligned}
 \tag{D.WI15}
\]
The indices in the second line are cyclic modulo $d$; those
used as independent coordinates run from zero to $d-2$.
Here $\mathcal Q_2$ is the ordinary square if $d$ is odd and
the Pontryagin square with values modulo $2d$ if $d$ is even.
Likewise $\mathcal Q_3$ is the ordinary cube if $3\nmid d$
and Thomas's Pontryagin third power with values modulo $3d$
if $3\mid d$.

Consequently, for every $d>0$ all the simply connected data
(D.WI1) give one oriented diffeomorphism type $M_d$, namely
the unique type with
\[
 H_2=\mathbb Z^{d-1}\oplus\mathbb Z/d,\qquad b_3=0,\qquad
 P_{\rm ref}=0,\qquad \mu\ \text{given by (D.WI15)}.
 \tag{D.WI16}
\]
One has $M_1=S^6$, and for $d$ prime to six this is the
type already identified in D.1.11. The types for different
$d$ are distinct, since their Euler numbers are $2d$.

\begin{proof}
\emph{Primitive component classes.}
Write $E$ for the complement of the interior of the cusp tube.
The actual finite-cap presentation before adjoining its cusp meridian is
\[
 H_1(E)=
 \langle\gamma,x_0',x_1'\mid m_0x_0'=n_0\gamma,\ 
                                   m_1x_1'=n_1\gamma\rangle .
\]
The determinant condition implies
$\gcd(m_0,m_1)=\gcd(m_i,n_i)=1$ and $H_1(E)=\mathbb Z$.
The cusp meridian is $x_0'+x_1'$ up to the common orientation sign.
It is primitive because adjoining its relation gives $H_1(X)=0$.
Angular differences of adjacent toric meridians are killed in
$H_1(E)$.  Thus the relative Thom map is
\[
 H^1(E)\longrightarrow H^2(X,E)=\mathbb Z^d,\qquad
                   1\longmapsto\pm(1,\ldots,1).
 \tag{D.WI4}
\]
The cokernel of the resulting injection
$\mathbb Z^d/\mathbb Z(1,\ldots,1)\hookrightarrow H^2(X)$
embeds in $H^2(E)$, which is torsion-free by the integral universal
coefficient theorem.  It has rank zero by D.2.34, so it vanishes.
In particular the entire fibre class is
$f_X^*[\mathrm{point}]=\sum_i x_i=0$.
For a reduced multiple fibre $D$ one has
$m[D]=f_X^*[\mathrm{point}]=0$; the group $H^2(X)$ is free,
so $[D]=0$ as well.

The common ruling in the $d$-period cusp is the specialization
of $a\wedge b$, with its complex orientation.  This is an integral
unit-period annulus calculation: that torus remains in one gamma
component and limits to the two adjacent external curves
$E_2+E_3=H-E_1$.  The three support functions
$r(r-1)/2$, $rs$, and $s(s-1)/2$ evaluate on that ruling as
$(-1,0,0)$, exactly their actual multiplier windings on
$a\wedge b$; its dual-torus coordinate is zero.
There is no division by $d$ or saturation at this step.
The primitive functional of D.2.34 evaluates it as
$\xi_d(a\wedge b)=1$.  The matrix (D.WE3) shows that $s$ has
exact order $d$ in the simply connected sector.  The actual
external intersections and the principal sum of components
therefore give (D.WI3).  The adjacent-pair formula counts both
edges when they have the same pair of labels.  Since every
product is torsion, its further product with $H^2(X)$ vanishes
in $H^6(X)=\mathbb Z$.

\emph{The differential and its literal critical cycle.}
The normalization immersion $j_i:dP_6\longrightarrow X$ has
normal line $K_{dP_6}$.  Indeed the principal local central
divisor is the product of the normal coordinates to all its
branches; on a fixed normalized branch its normal transition
is the inverse of the six toric boundary transitions.  This
includes both conductor branches of each self-identification.
Consequently
\[
 j_i^*c_1(TX)=c_1(TdP_6)+c_1(K_{dP_6})=0,\qquad
 \langle j_i^*c_2(TX),[dP_6]\rangle=6-6=0 .
 \tag{D.WI5}
\]
Restriction of $H^2(X)$ to the normalized components is injective.
For a component combination with coefficients $u_i$, its
restriction to component $i$ is
$\sum_r(u_{i+r}-u_i)C_r$.  Pairing with the anticanonical ample
class, which has degree one on each of the six boundary curves,
gives the component-graph Laplacian.  The graph is connected, so
zero restrictions force the $u_i$ to be constant.  That is
precisely the already removed principal relation.  Equation
(D.WI5) now gives $c_1(TX)=0$ integrally.

Remove the divisorial multiple-fibre factors from $df_X$.  Its
corrected target is the line
\[
 L=f_X^*T\mathbb P^1\otimes
                      \mathcal O\bigl(-\sum (m_i-1)D_i^{\rm fin}\bigr).
\]
It is topologically trivial by (D.WI4).  The abstract rank-two
kernel $V$ is locally free, including at $xy$ and $xyz$ charts,
and has the nowhere-zero abstract section given by $\delta$.
At a node the kernel columns are $(x,-y,0)$ and $(0,0,1)$;
at a triple point they are $(x,-y,0)$ and $(x,0,-z)$.
The constant scalar-circle weight vector is a nonzero constant
combination in these abstract bases.  At the good finite caps
the corrected differential is a submersion and $\delta$ remains
a nonzero tangent vector.  Since
$c_1(V)=c_1(TX)-c_1(L)=0$, its quotient by this section is a
trivial complex line.  Thus $V$ and $L$ admit full smooth frames.
They give the reference frame $\rho$ on the complement of the
literal critical curve cycle $C$.

The local Bott clutch at $xy$ has coefficient one:
its two normal columns are
$\bigl(\begin{smallmatrix}x&\overline y\\-y&\overline x\end{smallmatrix}\bigr)$.
The $xyz$ Hilbert--Burch minors give multiplicity one on each
of its three axes and no further codimension-two contribution.
The good finite logarithmic caps have no residual critical
cycle after correcting the target.  Hence, with the complex
orientation and the spin normalization of D.2.44--D.2.45,
\[
                 \lambda(TX,\rho)=-\operatorname{PD}(C).
 \tag{D.WI6}
\]

Only a local covering is used to compute $[C]$.
In the positive period basis
$N b=a$, $N\gamma_1=\delta$, the weight-one cusp's literal
three-curve sum has regular-fibre preimage
$a\wedge b+b\wedge\delta+\delta\wedge\gamma_1$.
The weight-$d$ cusp is its degree-$d$ etale period cover, with
$\gamma_d=d\gamma_1$.  Transfer therefore gives the actual
preimage of its full $3d$-curve sum as
\[
 d\,a\wedge b+d\,b\wedge\delta+\delta\wedge\gamma_d.
 \tag{D.WI7}
\]
The angular vector $\delta$ here is fixed by
$N\gamma_d=d\delta$; this convention removes the possible
opposite choice for the original scalar-circle parameter.
Since $\xi_d$ has values $1,0,-6d$ on these three bivectors,
(D.WI7) has image $-5d\,s=0$.  Thus $C$ is integrally
nullhomologous and $c_2(TX)=\lambda(TX)=0$.
No covering of either finite logarithmic cap is assumed.

\emph{A finite-cap detector with its actual reference parity.}
If $d$ is odd, $2P_{\rm ref}=\lambda=0$ already gives
$P_{\rm ref}=0$.  Suppose $d$ is even.  In the order-four
finite cap, put $\beta=[a]$ and $\eta=[u_1]$ in its two
central base directions.  Its marked base lattice is
\[
 \langle\beta,\eta,g\mid m_1g=-2r_1\beta+n_1\eta\rangle .
 \tag{D.WI8}
\]
Choose integers $B,C$ with $m_1C=-2r_1+n_1B$ and define the
primitive character by $\chi(\beta)=1$, $\chi(\eta)=B$,
$\chi(g)=C$.  Such integers exist because $\gcd(m_1,n_1)=1$.
Its connected fibre $N$ in the central four-manifold is an
elliptic-torus mapping torus.  In its base kernel
$h=g-C\beta$ and $q=\eta-B\beta$ obey $m_1h=n_1q$.
A primitive generator has the root-loop exponent inverse to
$n_1$ modulo $m_1$, hence odd modulo four.  The elliptic
monodromy of $N$ is therefore the order-four matrix or its
inverse.  Its Wang sequence gives
\[
             H_1(N;\mathbb Z)=\mathbb Z\oplus\mathbb Z/2.
 \tag{D.WI9}
\]
Moreover $\chi(\delta)=-2$, so the actual scalar tangent
vector is transverse to $N$ in the reduced finite fibre.
Projecting $J\delta,\operatorname{Re}v_2,\operatorname{Im}v_2$
to $TN$ along $\delta$ identifies the stable tangent frame
induced by $\rho$ with an actual tangent parallelization.
The normal three-frame is $\delta$ followed by the two real
directions of $L$.

For an actual parallelization of a closed oriented
three-manifold, its stable framed coefficient modulo two is
$1+\dim H_1(N;\mathbb F_2)$.  The spin-compatible Kaplan
zero-/two-handle proof was given in the finite-II frame lemma:
the relative $\lambda$ reduces to the relative Euler number,
and the even intersection matrix has even mod-two rank.
It applies to every prescribed spin structure
\cite[Definition 1.7 and Theorem 3.1, pp.239,243--246]{mapsKaplan79}.
Equation (D.WI9) therefore gives reference parity
\[
                         \epsilon=[N]_\rho=1\pmod2 .
 \tag{D.WI10}
\]

Retain the actual Thom chains, not only their matrix classes.
Let $K_{a0},K_{a1}$ be the first columns of the two blocks in
(D.WE3).  Directly in its four target coordinates,
\[
                    d s=-2\ell P\,K_{a0}-kP\,K_{a1}.
 \tag{D.WI11}
\]
The chain at cap zero misses $N$, whereas the $a$-Thom
chain in cap one pairs with $N$ by $\chi(a)=1$.
Thus $\operatorname{lk}(N,s)=-kP/d$.
The determinant equation implies that $k$ is odd, so $N$
generates the two-primary summand of $H_3(X)$.
The specialization chains for (D.WI7) are supported in the
cusp and miss $N$.  Equations (D.WI7)--(D.WI11) then give
the actual integer link
\[
                          \ell(N,C)=5kP,
 \tag{D.WI12}
\]
up to the common harmless orientation sign.  It is odd.

Finally the relative refined-class computation extends with
reference parity $\epsilon$ in the form
\[
 \operatorname{lk}(N,\operatorname{PD}P_{\rm ref})
                 =-\frac{\ell(N,C)}2-\frac{\epsilon}2
                                      \quad\text{in }\mathbb Q/\mathbb Z .
 \tag{D.WI13}
\]
Indeed, choose Zhubr's admissible three-frame on a genuinely
embedded torsion sphere obtained from an odd multiple of $N$
by spin-compatible ambient two-handle surgery.  Its complement
spin three-frame extends over the trace.  The induced stable
framed class on $N$ is even.  The comparison with $\rho$ has
Spin(6) degree $q_0\equiv\epsilon\pmod2$.
For a bounding four-chain $B_4$ with boundary
$\operatorname{ord}(N)N$, the exact relative transgression is
$-B_4\cdot C=2q+\operatorname{ord}(N)q_0$.
Division of this integer identity gives (D.WI13), retaining
the actual rational link before reduction modulo one.
This is the proof of the relative lemma in D.2.45 with its
reference-parity term restored, not an application of the
unproved general change-of-frame conjecture.
Free $H_2$ causes no change: $C$ is nullhomologous, and
$2P_{\rm ref}=\lambda=0$ already kills the free part.

The two odd halves in (D.WI13) cancel by (D.WI10) and
(D.WI12).  The two-primary part of $P_{\rm ref}$ is zero,
and its odd-primary part was already zero from
$2P_{\rm ref}=0$.  This proves (D.WI2).
\end{proof}

% FULL OPERATIVE CONTINUATION: sources/returns/S6_04_MAPS_MANIFOLDS/code/maps_weighted_engel_mu_proof.tex
% Continuation of D.2.48: the full invariant, with both exceptional moduli.
\begin{proof}[Proof of the complete invariant and classification]
\emph{A global finite-coefficient class, retaining boundary torsion.}
At weight one the regular bundle $V_1$ over the peripheral
graph has free $H_1$ and $H_2$. Its degree-two group has the
primitive fibre quotient generated by $s$, together with the
five marked graph generators in (D.WE11). Choose the integral
cohomology class $\alpha_1$ which evaluates as $\xi_1$ on the
fibre part and zero on every graph generator.

On the weight-one cusp use the actual toric line with support
\[
                 \psi(r,s)=-\frac{r(r-1)}2+3s(s-1).
 \tag{D.WI17}
\]
Its multiplier windings give $\xi_1$ in the positive period
marking used in (D.WI7). Its boundary value agrees integrally
with $\alpha_1$: their fibre restrictions agree, and both
vanish on the two angular Wang sweeps, which bound in the
cusp. These are all the remaining generators of $H_2(M_1)$.
Also $H_1(M_1)$ is free, so the integral universal coefficient
theorem leaves no hidden $\operatorname{Ext}$ difference.

Before the finite caps are attached, the core and cusp have
the genuine degree-$d$ period covering defined by
$\gamma_d=d\gamma_1$. Pull back both classes along that
cover. They then agree on the entire boundary $M_d$, not
only on its regular-fibre restriction. In particular this
argument retains the $\mathbb Z/d$ term in $H_1(M_d)$ and
the resulting torsion in $H^2(M_d)$. The core pullback
$\alpha_d$ has fibre functional $\xi_d$ and zero values on
the five graph generators: the covering fixes the first
three such generators and multiplies the two $u_i$ generators
by $d$.

At a finite cap let $\pi:Y_i\longrightarrow Q_i$ be its
actual normal circle bundle. The obstruction for
$\alpha_d\bmod d$ to extend across the disk bundle is its
Gysin image
\[
                  \pi_!(\alpha_d\bmod d)\in
                         H^1(Q_i;\mathbb Z/d).
 \tag{D.WI18}
\]
Its values on the complete marked generators
$a,\gamma,g$ of $H_1(Q_i)$ are precisely the first-row
coefficients of (D.WE3):
\[
 \begin{array}{c|ccc}
       &a&\gamma&g\\ \hline
 i=0&-2dn_0&-6dr_0&0\\
 i=1&-3dn_1&-6dr_1&0.
 \end{array}
 \tag{D.WI19}
\]
The root-loop column $g$ is retained. This is the Thom
definition of fibre integration: a loop in $Q_i$ sweeps its
normal circle, and pairing with $\alpha_d$ gives the displayed
actual normal-column coefficient. Every coefficient is zero
modulo $d$. Since $H^1(Q_i;\mathbb Z/d)$ is
$\operatorname{Hom}(H_1(Q_i),\mathbb Z/d)$, (D.WI18) vanishes,
even when the integral normal line has torsion. The circle
Gysin sequence therefore supplies an extension over each
finite cap. Mayer--Vietoris now glues them with the pulled-back
cusp line to a class $u\in H^2(X;\mathbb Z/d)$.
Only the core and cusp were covered; neither finite fill was
assumed to be a cover of a weight-one fill.

On every normalized cusp component the integral representative
of its restriction is
\[
              \eta=-5H+6E_1-E_3,\qquad
       \eta(E_1,E_2,E_3)=(-6,0,1),\qquad
       \eta(f)=1,\qquad \eta^2=-12.
 \tag{D.WI20}
\]
The three degrees follow directly by taking the differences
of (D.WI17) across the internal, diagonal, and scalar-fixed
walls, as in D.1.9. In particular $u(s)=1$.
Choose free lifts in $H_2(X)$ annihilated by $u$, subtracting
appropriate multiples of $s$. This specifies the splitting
in (D.WI14)--(D.WI16). The integral classes $x_i$ form its
dual free basis, and $u$ is the standard cyclic coordinate.

\emph{The full mixed coordinates.}
The first mixed value in (D.WI15) follows from (D.WI3) and
$\langle u\kappa,[X]\rangle=1$. The Pontryagin square of the
reduction of an integral degree-two class is its ordinary
square with the full target modulus $2d$; this is the
$p=2$ case of the actual cochain construction
\cite[equations (2.1), (2.8), and (3.4), pp.267--268]{ffThomas56}.
By naturality and the projection formula for the normalized
component fundamental class,
\[
 \langle \mathcal Q_2(u)x_i,[X]\rangle
       =\langle\mathcal Q_2(j_i^*u),[dP_6]\rangle
       =\eta^2=-12\pmod {D_2}.
 \tag{D.WI21}
\]
This computes the coefficient modulo $2d$ when $d$ is even;
reducing the ordinary square merely modulo $d$ would lose
precisely the extra order-two information. The relation
$\sum_i x_i=0$ is consistent with (D.WI21), since
$12d$ is divisible by $D_2$.

\emph{The pure coordinate.}
The following actual scalar-circle calculation evaluates
the full Pontryagin third power. It uses the class just
constructed and its specified component restrictions.
\input{sections/fourfolds_D48_thomas.tex}

\emph{Why these detectors are complete.}
We include the low-degree integral target calculation to
exclude an undetected finite subgroup. For a cyclic group,
\[
 \begin{aligned}
 H_2(K(\mathbb Z/d,2);\mathbb Z)&=\mathbb Z/d,\\
 H_4(K(\mathbb Z/d,2);\mathbb Z)&=\mathbb Z/D_2,\\
 H_6(K(\mathbb Z/d,2);\mathbb Z)&=\mathbb Z/D_3 .
 \end{aligned}
 \tag{D.WI22}
\]
The degree-four assertion follows from Whitehead's exact
sequence: on $K(A,2)$ its two adjacent homotopy groups
vanish, so $H_4=\Gamma(A)$
\cite[Theorem 1 and Section 2]{mapsWhitehead50}.
For $A=\mathbb Z/d$, the universal quadratic generator
$q(1)$ has relations $d^2q(1)=2dq(1)=0$.
Conversely $q(j)=j^2$ modulo $\gcd(d^2,2d)=D_2$
is well defined and quadratic. Thus these are its complete
relations and the image of $\mathbb {CP}^2$ generates.
The ordinary or Pontryagin square evaluates on that image
as one modulo $D_2$, so it detects the entire group.

For degree six, Zhubr's epimorphism in Section 3.11 and
his two free bordism generators in Section 3.12 show that
the spin-realizable part for a cyclic coefficient group
is generated in homology by the image of
$[\mathbb {CP}^3]$: his other generator has zero
fundamental-class image. For odd-primary cyclic groups
this is all of $H_6$, since the only obstruction in his
Lemma 3.7 has mod-two target. For a two-power cyclic
group Section 3.12 explicitly gives
$\ker\alpha=H_6=\mathbb Z/2^a$
\cite[Lemmas 3.7 and 3.11, and Section 3.12,
pp.805--807]{manZhubr75}.
The decomposition of $\mathbb Z/d$ into its relatively
prime primary factors gives the corresponding product
of Eilenberg--MacLane spaces. All cross tensor and
$\operatorname{Tor}$ terms vanish, since their positive
homology has relatively prime orders. Hence $H_6$ is
cyclic and is generated by the image $z$ of
$[\mathbb {CP}^3]$ for every $d$.

The order of $z$ divides both $d^3$ and $3d$.
Indeed the class $d h$ on $\mathbb {CP}^3$ becomes zero
modulo $d$, while its integral cube is $d^3$.
On $\mathbb {CP}^1\times\mathbb {CP}^2$ the class
$d x+y$ reduces to a class factoring through
$\mathbb {CP}^2$, whose degree-six image is zero,
while its integral cube evaluates to $3d$.
Thus the order divides
$\gcd(d^3,3d)=d\gcd(d,3)=D_3$.
The opposite bound is detected by the ordinary cube when
$3\nmid d$, or by Thomas's operation when $3\mid d$:
on the reduced hyperplane class of $\mathbb {CP}^3$
the full target value is one
\cite[equation (2.8), p.267]{ffThomas56}.
This proves the last assertion of (D.WI22) with its
stated perfect detector.

Finally
$K(\mathbb Z^r\oplus\mathbb Z/d,2)
 =(\mathbb {CP}^{\infty})^r\times K(\mathbb Z/d,2)$.
The first factor has torsion-free homology in even
degrees and zero homology in odd degrees. The integral
Kunneth formula has no $\operatorname{Tor}$ term and
gives exactly (D.WI14). Integral cubic monomials,
$u x_i x_j$, $\mathcal Q_2(u)x_i$, and
$\mathcal Q_3(u)$ are the respective perfect coordinate
detectors. Therefore (D.WI15) computes all of $\mu$,
including the extra factors two and three.

\emph{Application of smooth classification.}
The manifold is simply connected and spin, has
$H_2=\mathbb Z^{d-1}\oplus\mathbb Z/d$ and middle Betti
number zero by D.2.34, refined class zero by the first
proof, and the complete fundamental-class invariant
(D.WI15). These are exactly the four invariants in
Zhubr's Theorem 1, with the definitions on pp.794--795
\cite[Theorem 1]{manZhubr75}. Their values are independent
of every marked parameter in (D.WI1).
Existence is supplied by the constructed manifolds,
and the injective classification gives one oriented
diffeomorphism type for each $d$. This proves (D.WI16)
by an applied complete classification, rather than by
ordinary cohomology or an asserted connected-sum model.
\end{proof}

% FULL OPERATIVE CONTINUATION: sources/returns/S6_04_MAPS_MANIFOLDS/sections/fourfolds_D48_thomas.tex
% Primary-operation input for maps-owned D.2.48; no separate theorem number.
\subsubsection*{The full Pontryagin cube of the weighted Engel torsion class}
We use the actual class $u\in H^2(X_d;\mathbb Z/d)$ whose
restrictions to the normalized cusp components are specified
in D.2.48. In particular its degree on every scalar-fixed
rational curve $C_j$, $j\in\mathbb Z/d$, is one, and its
component restriction is represented by the integral toric
support
\[
                  \psi(r,s)=-\frac{r(r-1)}2+3s(s-1).
                                                               \tag{D.TP1}
\]
The calculation below uses these actual restrictions. Merely
choosing an arbitrary degree-two class that is a unit on
torsion would not fix its pure coordinate in the presence
of the free summand.

The scalar circle action is semifree. On a finite cap of
order $m_i$ its gamma translation has primitive numerator
$n_i$, with $\gcd(m_i,n_i)=1$. A nontrivial root power cannot
be compensated by the scalar direction, because that would
require $j n_i/m_i$ integral. Thus these caps are scalar-free.
The cusp has exactly the $d$ fixed curves $C_j$; their normal
bundles are $\mathcal O(-1)\oplus\mathcal O(-1)$ with weights
$(1,-1)$, and the other cusp orbits are free. These are the
marked geometric inputs of the weighted family.

Because $X_d$ is simply connected, $u$ has an equivariant
lift $\widehat u\in H^2_{S^1}(X_d;\mathbb Z/d)$.
In the Borel spectral sequence the only possible
differentials out of ordinary degree two have targets
involving $H^1(X_d;\mathbb Z/d)$ or $H^3(BS^1;\mathbb Z/d)$,
and both vanish. Write
\[
         \widehat u|_{C_j}=x+\lambda_j t,
         \qquad x=[C_j]^*,\quad t\in H^2(BS^1;\mathbb Z/d).
                                                               \tag{D.TP2}
\]
The difference of any two equivariant lifts of a fixed
ordinary class on the simply connected normalized
$dP_6$ is a multiple of $t$. Hence the differences of
the characters on its two scalar-fixed curves can be
read from any toric equivariant lift of (D.TP1).
The fixed edges are the translates of the edge from
$(0,j)$ to $(0,j+1)$; the support slope in their scalar
direction is
\[
               \psi(0,j+1)-\psi(0,j)=6j.
                                                               \tag{D.TP3}
\]
The fibre character of the associated monomial
trivialization is the negative of this slope, up to
one global character convention. Equivalently the
intervening weight-one orbit curve has degree $-6$.
Thus, for one common sign $\epsilon=\pm1$,
\[
                    \lambda_j=\lambda_0+6\epsilon j
                                 \quad\text{in }\mathbb Z/d.
                                                               \tag{D.TP4}
\]
This equation is compatible with the period $j\mapsto j+d$.
It does not assert that the integral support line itself
descends with an integral global circle linearization:
the needed character equality is modulo $d$ and follows
from the existing global equivariant cohomology lift.

Suppose now that $3\mid d$. Thomas's operation is defined
with precisely the required full coefficients:
\[
 \mathcal P_3:H^2(-;\mathbb Z/d)
                    \longrightarrow H^6(-;\mathbb Z/(3d)).
                                                               \tag{D.TP5}
\]
This is \cite[equation (2.1), p.267]{ffThomas56} with
$p=3$ and $m=d/3$. On the reduction of an integral
class $z$, it equals the reduction of $z^3$ modulo
$3d$. Indeed the correction term in Thomas's actual
cochain formula (3.4) is zero for an integral cocycle,
and its first term is the ordinary threefold diagonal
product \cite[equations (3.4)--(3.6), p.268]{ffThomas56}.
Each restriction (D.TP2) has such an integral lift.

Apply equivariant localization with coefficients
$\mathbb Z/(3d)$ to $\mathcal P_3(\widehat u)$.
No field assumption is needed here: the complement
is scalar-free, and each normal Euler class becomes
a unit after inverting $t$. On a fixed curve it is
\[
                    (-x+t)(-x-t)=-t^2,
\]
because $x^2=0$. Integrating the term of ordinary
curve degree two gives
\[
 \left\langle\mathcal P_3(u),[X_d]\right\rangle
                      =-3\sum_{j=0}^{d-1}\lambda_j^2
                                  \quad\text{in }\mathbb Z/(3d).
                                                               \tag{D.TP6}
\]
The right side is independent of the integer
representatives of the characters: $3\mid d$ makes
the cube of an integral lift well defined modulo
$3d$. The base cohomology polynomial ring injects
into its localization at $t$, so the degree-zero
identity is an identity in the original coefficient
ring, not just a localized statement.

For arbitrary integers $d>0,\lambda_0$ and either
sign in (D.TP4), the exact identity is
\[
 \sum_{j=0}^{d-1}(\lambda_0+6\epsilon j)^2
   =d\left(\lambda_0^2+6\epsilon\lambda_0(d-1)
                         +6(d-1)(2d-1)\right).              \tag{D.TP7}
\]
Substitution into (D.TP6) proves
\[
                 \left\langle\mathcal P_3(u),[X_d]\right\rangle=0
                                  \quad\text{in }\mathbb Z/(3d).
                                                               \tag{D.TP8}
\]
Thus the complete pure cyclic fundamental-class coordinate
is zero, including the order-three subgroup invisible to
the ordinary mod-$d$ cube. For $3\nmid d$ the same
localization calculation with the ordinary cube gives
zero modulo $d$. Mixed coordinates are separate from
this calculation; in particular it does not discard
the order-$2d$ mixed Pontryagin-square terms when $d$
is even.

The universal integer identities, support slope,
single-curve Euler/cube calculation and an explicitly
bounded supplementary check are executed in
\path{code/fourfolds_D48_thomas.py}.
The supplementary range is $d=3,6,\ldots,120$,
$\lambda_0=0,\ldots,d-1$ and both signs, giving
$4920$ exact checks. The arbitrary-$d$ conclusion
comes from (D.TP7), not that finite range.
